\section{Related Work}
\label{sec:related-work}

\subsection{World Foundation Models for Physical AI}

The concept of learning internal models of the environment, commonly known as world models, has a long history in reinforcement learning and cognitive science~\cite{ha2018world}. Early approaches such as Dreamer~\cite{hafner2020dreamer} and DayDreamer~\cite{wu2023daydreamer} learn compact latent-state dynamics models that enable planning through imagined rollouts, demonstrating the viability of model-based learning for physical robot control. More recently, the success of large-scale generative models has motivated a paradigm shift toward \emph{world foundation models} (WFMs) that operate directly in pixel or video space. Sora~\cite{liu2024sora} demonstrated that video diffusion models can generate temporally coherent, physically plausible visual simulations. The Cosmos platform~\cite{nvidia2025cosmos} further formalized this direction by training DiT-based~\cite{peebles2023dit} video generation models on internet-scale data for physical AI applications, while Cosmos 3~\cite{nvidia2026cosmos3} extended it to an omnimodal framework unifying language, image, video, audio, and action modalities. CogVideoX~\cite{yang2024cogvideox}, HunyuanVideo~\cite{kong2024hunyuanvideo}, Stable Video Diffusion~\cite{blattmann2023svd}, and Wan~\cite{wan2025wan} achieved high-fidelity video generation across diverse domains. DIAMOND~\cite{alonso2024diamond} demonstrated that diffusion-based world models can capture fine visual details critical for decision-making. Pelican-Unified~\cite{zhang2026pelican} advocates a unified paradigm integrating world modeling with understanding, reasoning, and action for embodied intelligence. However, these models are fundamentally single-view: they generate one coherent video stream without explicit mechanisms for multi-view consistency.

Interactive world models represent another important direction. UniSim~\cite{yang2023unisim} learns real-world simulators from diverse data sources, Genie~\cite{bruce2024genie} and Genie 2~\cite{parkerholder2024genie2} enable interactive environment generation from single images, and iVideoGPT~\cite{wu2024ivideogpt} scales autoregressive world models to complex environments. In the robotic manipulation domain, GR-2~\cite{cheang2024gr2} builds a generative video-language-action model with web-scale knowledge, IRASim~\cite{zhu2025irasim} learns action-conditioned video simulators from real robot data, EnerVerse~\cite{huang2025enerverse} envisions embodied future spaces for manipulation planning, and LaDi-WM~\cite{huang2025ladiwm} employs latent diffusion-based world models for predictive manipulation. AgiBot World~\cite{bu2025agibotworld} provides a large-scale multi-view manipulation platform, and WorldArena~\cite{shang2026worldarena} offers a unified benchmark for evaluating embodied world models. While some of these systems support multiple viewpoints via token concatenation, none introduces explicit cross-view geometric reasoning---a critical limitation for robotic applications, where 3D consistency directly governs policy quality.

\subsection{Multi-View and 3D-Aware Visual Generation}

Multi-view generation has been extensively studied in the context of 3D content creation from single images. Zero-1-to-3~\cite{liu2023zero123} pioneered viewpoint-conditioned diffusion by fine-tuning a 2D diffusion model with relative camera transformations. SyncDreamer~\cite{liu2023syncdreamer} introduced synchronized multi-view denoising with 3D-aware attention to ensure geometric consistency across generated views. MVDream~\cite{shi2023mvdream} extended multi-view diffusion to text-conditioned 3D generation, while MVDiffusion~\cite{tang2023mvdiffusion} enabled holistic multi-view image generation with correspondence-aware attention. SV3D~\cite{voleti2024sv3d} leveraged latent video diffusion for orbital multi-view synthesis, and SV4D~\cite{xie2024sv4d} further extended this to dynamic 3D content with multi-frame and multi-view consistency. CAT3D~\cite{gao2024cat3d} demonstrated that multi-view diffusion models can create high-quality 3D assets from arbitrary inputs. These methods demonstrate that explicit view-conditioning and cross-view interaction mechanisms can produce 3D-consistent multi-view images. However, they are designed for \emph{static} or short-sequence object-centric scenes and do not address generating consistent multi-view \emph{videos} of dynamic environments as required by world models.

Camera control in video generation has been explored by CameraCtrl~\cite{he2024cameractrl}, which represents camera poses using Pl\"ucker ray coordinates and injects them into video diffusion models. ViewCrafter~\cite{yu2024viewcrafter} further tames video diffusion models for high-fidelity novel view synthesis. These works provide technical foundations for camera-aware generation, but focus on single-view camera trajectory control rather than multi-view consistency.

\subsection{3D Representations and Geometric Reconstruction}

Recent advances in geometric 3D vision have produced powerful tools for evaluating and enforcing 3D consistency. Neural Radiance Fields (NeRF)~\cite{mildenhall2020nerf} and 3D Gaussian Splatting~\cite{kerbl2023gaussiansplatting} established the foundation for differentiable 3D scene representations. DUSt3R~\cite{wang2024dust3r} and its successor MASt3R~\cite{leroy2024mast3r} learn to predict dense 3D point maps from image pairs without requiring known camera parameters, enabling direct measurement of cross-view geometric consistency. The Depth Anything series~\cite{yang2024depthanything,yang2024depthanythingv2,lin2025depthanything3} provides foundation models for monocular depth estimation that capture robust 3D-aware features from large-scale training. VGGT~\cite{wang2025vggt} further advances this direction by grounding visual geometry in a unified transformer that jointly predicts camera poses, depth, and 3D point maps from arbitrary image collections. These models encode rich geometric priors that can serve as supervision targets for world models seeking 3D consistency.

The REPA framework~\cite{yu2024repa} introduced the idea of aligning intermediate representations of diffusion transformers with those of pretrained encoders, demonstrating accelerated training and improved generation quality for image diffusion. Our Latent 3D-REPA extends this principle from 2D image generation to multi-view video world models, using 3D-aware features from VGGT and Depth Anything 3 as the alignment target to inject geometric consistency directly into the diffusion process.

% \subsection{Robotic World Models and Embodied AI}

% The application of world models to robotics has gained significant momentum. RoboDreamer~\cite{zhou2024robodreamer} learns compositional world models that enable robot imagination for planning. The Open X-Embodiment initiative~\cite{collaboration2023openx} and associated RT-X models provide large-scale multi-robot datasets and policies, establishing the standard multi-camera configuration (wrist, overhead, egocentric) for robotic manipulation. RT-1~\cite{brohan2023rt1} and RT-2~\cite{brohan2023rt2} demonstrated scalable robotic control via transformers, while BridgeData V2~\cite{walke2023bridgedata} and RoboCasa~\cite{nasiriany2024robocasa} provide diverse manipulation data. Octo~\cite{octo2024} offers a generalist policy trained on multi-embodiment data, and $\pi_0$~\cite{black2024pi0} shows that vision-language-action flow models can achieve general robot control. NVIDIA's GR00T~\cite{nvidia2025groot} further scales foundation models for humanoid robots.

% These works highlight that modern robotic systems fundamentally rely on multi-view observations, yet the world models that serve them remain predominantly single-view. \methodname{} bridges this gap by providing a 3D-consistent multi-view world model backbone that can directly serve as a simulator for multi-camera robotic planning and policy learning.

% \subsection{Diffusion Transformers and Positional Encoding}

% Our work builds upon the Diffusion Transformer (DiT) architecture~\cite{peebles2023dit}, which replaces U-Net backbones with transformer blocks for diffusion-based generation. DiT's flexibility in handling variable-length token sequences makes it naturally suited to multi-view inputs. Rotary Position Embedding (RoPE)~\cite{su2024rope} provides an efficient mechanism for encoding positional information that preserves relative position relationships, a property we exploit in our Geometric Rotary Position Embedding to encode camera geometry directly into the attention mechanism. The latent diffusion paradigm~\cite{rombach2022ldm,ho2020ddpm} operates in a compressed representation space, which both reduces computational cost and provides the natural setting for our latent-level representation alignment.
